\documentclass[11pt]{article}
\usepackage{wsi-style-lua}
\usepackage{amsmath}
\newcommand{\Nesc}{\mathcal{N}_{\mathrm{esc}}}
\wsirunhead{Spacetime as Escrow Bookkeeping}
\title{Spacetime as Escrow Bookkeeping:\\[2pt]
\large A Conceptual Translation of General Relativity into the\\ Static Entropy Escrow Vocabulary}
\author{Grant Whitmer}
\date{May 2026}
\begin{document}
\maketitle
\begin{abstract}
We recast four standard results of general relativity---gravitational time dilation, the Tolman temperature law, the Bekenstein--Hawking entropy, and Jacobson's (1995) thermodynamic derivation of the Einstein equations---in the vocabulary of the static gravitational entropy escrow framework [1]. A single ratio, $S_{\mathrm{esc}} = |U_{\mathrm{grav}}|/T_U$ (gravitational binding energy over local Unruh temperature), reproduces all four: it equals the Bekenstein--Hawking entropy of a Schwarzschild horizon exactly, supplies the entropy flux in Jacobson's first law, and governs the Tolman redshift. No underlying physics is modified; the contribution is interpretive---a thermodynamic re-reading, not a new equation of motion. We are explicit about the limits: $|U_{\mathrm{grav}}|$ takes regime-specific forms not yet unified by a single covariant construction, and $T_U$ appears in two conventions related by the Tolman shift, so the unification is exact at the algebraic level $S = |U|/T$ but only partial as a single observable. Lattice tests [2] recover the Bisognano--Wichmann asymptote partially in $1{+}1$D (with an unexplained $1/30$ prefactor) and show a qualitatively different form in $3{+}1$D. Deriving that prefactor from first principles is the open calculation that would turn this conceptual synthesis into a constraining proposal.
\end{abstract}

\section{Introduction}


\subsection{What this paper is}

The static gravitational entropy escrow framework [1] proposes that
gravitational binding energy is entropy held "in escrow" against the
local Unruh temperature: for any gravitationally bound configuration
with binding energy $|U_{\mathrm{grav}}|$ at a location where the local Unruh
temperature is $T_{U}$, the escrow entropy is $S_{\mathrm{esc}}$ = $|U_{\mathrm{grav}}|$/$T_{U}$. The
framework paper develops the consequences of this single relation,
including the recovery of Newton's law via an entropic-force argument
and a reproduction of the Bekenstein--Hawking entropy formula at
horizons. Companion work [2] tests the modular-Hamiltonian
interpretation of the postulate against free scalar field theory on
regulated lattices in 1+1 and 3+1 dimensions; that work reports
approximate recovery of the Bisognano--Wichmann linear asymptote in 1+1D
at small distances from the cut, with a prefactor suppression of
approximately 1/30 of the literal BW value, and non-recovery in 3+1D
within the d1 range the present lattice can resolve.

What the framework paper does not do, and what the lattice paper does
not do, is connect the escrow postulate to the established literature on
the thermodynamic derivation of Einstein's equations --- specifically
to Jacobson's (1995) derivation [3] from $\delta Q$ = T·dS at local Rindler
horizons, to the Tolman--Ehrenfest relation [4] for thermal
equilibrium in static gravitational potentials, and to the
Faulkner--Guica--Hartman--Myers--Van Raamsdonk (FGHMV) result [5,6]
that demanding the entanglement-entropy first law for vacuum-state
perturbations gives linearized Einstein equations in asymptotically AdS
spacetimes. These three results, plus the Bekenstein--Hawking formula
[7,8], constitute a now-substantial body of work establishing that
general relativity has the standing of a thermodynamic equation of
state. The static escrow postulate, as stated, lives squarely inside
this body of work. The connection has not previously been made explicit.

The present paper makes the connection explicit. We translate each of
these four standard results into the escrow vocabulary, show that the
translation in each case is mathematically empty (no equation of
standard GR is modified or supplemented), and demonstrate that the
escrow postulate provides a single interpretive object --- $S_{\mathrm{esc}}$ ---
that lets the four results be read as four faces of one identity. We are
explicit (as developed in §V.G--H) that this 'single object'
description is partly notational: the $|U_{\mathrm{grav}}|$ feeding into $S_{\mathrm{esc}}$
takes regime-specific forms across the four translation legs, and $T_{U}$ is
used with two related but distinct conventions; the unification is real
at the level of the algebraic form S = |U|/T, and partial at the level
of identifying a single covariant observable. The contribution is a
conceptual re-reading, not new physics. We are explicit throughout that
'escrow' is interpretive vocabulary, not a claim about a literal
account-keeping mechanism in the vacuum. $S_{\mathrm{esc}}$ = $|U_{\mathrm{grav}}|$/$T_{U}$ is a
defined ratio; the bookkeeping language (held in escrow, released, flow)
is suggestive shorthand for the thermodynamic structure of that ratio,
not a claim about ontologically separate accounts or ledgers. The
framework's scientific content is the dimensional and structural
relation $S_{\mathrm{esc}}$ = $|U_{\mathrm{grav}}|$/$T_{U}$ and its consequences, not the linguistic
frame in which we describe it.


\subsection{Why a translation paper is worth writing}

Three observations motivate the present paper. First, the escrow
framework's published presentation [1] does not state its
relationship to Jacobson's derivation explicitly, with the consequence
that readers familiar with [3] have no easy path to placing the escrow
postulate within the existing thermodynamic-gravity literature. The
framework is at risk of being read as a parallel proposal competing with
[3] when it is, more accurately, a specification of what the S in
Jacobson's $\delta Q$ = T·dS is the entropy of. The present paper makes this
specification explicit and addresses the parallel-proposal misreading.

Second, the lattice paper [2] tests the modular-Hamiltonian leg of
this connection numerically and finds partial recovery in 1+1D and
non-recovery in 3+1D. Without an explicit statement of what is being
tested, readers of [2] may be unclear which interpretation of the
escrow postulate the lattice results actually constrain. The present
paper supplies the explicit interpretation that [2]'s tests target.

Third, the framework's history includes one documented failed extension
--- the C8 entropy-current attempt described in [9] --- in which an
attempt to make the postulate dynamically active reduced to a
rate-equation reformulation of the saturated Bekenstein bound. The
failure mode there was an overreach beyond what the framework's
interpretive content actually supports. The present paper consolidates
what the framework can and cannot claim at the interpretive level, with
explicit guardrails (§V) against the failure modes documented in [9]
and in the present authors' subsequent reviewer-feedback iterations on
[2]. These guardrails are the principal pedagogical contribution of
the paper, and they include (i) explicit rejection of a separate
"escrow stress-energy tensor" added to the matter source of
Einstein's equation, (ii) explicit rejection of any "vacuum
computational bandwidth" reading of the framework, and (iii) explicit
observer-dependence of the escrow entropy (different observers,
different escrows).


\subsection{Scope and reading order}

Sections II, III, and IV present the three legs of the translation: time
dilation as Tolman temperature shift (§II), the Jacobson derivation as
an escrow-flow identity (§III), and the modular-Hamiltonian connection
to vacuum entanglement structure (§IV). Each section follows a common
template: (A) standard statement of a known GR result, (B) translation
into escrow vocabulary, (C) why the two readings reduce to the same
identity, (D) explicit scope (what the section does and does not
establish), (E) explicit guardrails against the local failure modes
specific to that section. Section V consolidates the guardrails from
§II.E, §III.F, and §IV.E into a single statement of what the framework
does not commit to, with subsections G and H explicitly naming the
framework's unresolved scope-limitations (§V.G: $|U_{\mathrm{grav}}|$ lacks a
single covariant definition across regimes; §V.H: $T_{U}$ is used with two
related but distinct conventions related by the Tolman shift). Section
VI states the conclusion and identifies the principal open theoretical
problems.

Three contributions of this paper are worth flagging at the outset.
First, equation (8) of §II isolates the dimensionless geometric quantity
2πr/$\lambda_{C}$ --- the ratio of a circumference to the test particle's reduced
Compton wavelength --- as the dimensionless content of the escrow
entropy for a single test mass; we introduce the working notation $\mathcal{N}_{\mathrm{esc}}$
≡ 2πr/$\lambda_{C}$ for this quantity in §II.B and identify it as the candidate
organizing variable for the test-mass leg of the translation. This ratio
is the unique dimensionless invariant constructible from the
Newtonian-limit ingredients (binding energy GMm/r divided by an
acceleration-based Unruh temperature), so its appearance in equation (8)
is dimensionally forced rather than a discovery; what is non-trivial is
that the same dimensionless quantity recurs in independently-derived
corners of semiclassical gravity (saturation of the Bekenstein bound for
a test mass of energy mc² confined to radius r, Euclidean horizon
periodicity, modular flow generators), and the recurrence is the
smallest piece of evidence the framework is tracking a real organizing
variable rather than performing a sophisticated change of variables.
Second, equations (17)--(18) of §III show that the escrow postulate's
prediction for the entropy of a Schwarzschild black hole equals the
Bekenstein--Hawking entropy to all displayed digits, with no fudge
factor; we describe this as a structural-consistency check rather than
an independent prediction (the equality follows algebraically from the
choice $|U_{\mathrm{grav}}|$\textsuperscript{Schw} = Mc²/2 and the standard Hawking temperature
formula), and we extend the consistency check in §III.D to
Reissner--Nordström and Kerr horizons via the Smarr formula. Third,
§IV's connection to the FGHMV result places the lattice paper's
empirical content [2] inside the FGHMV theoretical framework: the 1/30
prefactor and the 3+1D non-recovery become specific calculational
questions about how lattice-regulated free QFT approaches its continuum
BW limit, rather than free-floating empirical curiosities.

The framework's broader empirical anchors live in [1] and depend on
the static escrow postulate's horizon-thermodynamic content, which is
unchanged by the conceptual reframing the present paper provides. We do
not revisit those anchors here; the present paper concerns the
GR-to-escrow translation.


\section{Time Dilation As Tolman Temperature Shift}


\subsection{The standard statement}

A clock at radius r outside a non-rotating mass M runs slow relative to
a clock at infinity by the Schwarzschild redshift factor

\begin{equation}
  \frac{d\tau}{dt} = \sqrt{1 - \frac{2GM}{rc^{2}}}
  \label{eq:1}
\end{equation}

where τ is proper time at r and t is coordinate time (the proper time of
an observer at infinity). In the weak-field limit GM/rc² ≪ 1,

\begin{equation}
  \frac{d\tau}{dt} \approx 1 - \frac{GM}{rc^{2}}
  \label{eq:2}
\end{equation}

This is the standard expression for gravitational time dilation.

A separate but related fact about static gravitational potentials is the
Tolman relation [4]. A thermal bath in equilibrium across a static
spacetime does not have a uniform local temperature; the local
temperature $T_{\mathrm{local}}$(r) measured by a stationary observer at radius r is
related to the temperature measured at infinity $T_{\infty}$ by

\begin{equation}
  T_{\mathrm{local}}(r)\,\sqrt{-g_{tt}(r)} = T_{\infty} = \text{constant}
  \label{eq:3}
\end{equation}

For the Schwarzschild metric this is

\begin{equation}
  T_{\mathrm{local}}(r)\,\sqrt{1 - \tfrac{2GM}{rc^{2}}} = T_{\infty}
  \label{eq:4}
\end{equation}

Comparing equations (1) and (4) shows that the local temperature is
boosted relative to the temperature at infinity by exactly the same
factor by which the local clock runs slow:

\begin{equation}
  \frac{T_{\mathrm{local}}(r)}{T_{\infty}} = \frac{dt}{d\tau} = \frac{1}{\sqrt{1 - 2GM/rc^{2}}}
  \label{eq:5}
\end{equation}

This is a known result. It says that gravitational time dilation and the
Tolman temperature shift are the same equation, expressed in two
different vocabularies. What the standard derivation does not provide is
an interpretation of why these two should be the same equation. The
escrow framework provides one.


\subsection{The escrow translation}

The static escrow postulate [1] asserts that for any gravitationally
bound configuration with binding energy $|U_{\mathrm{grav}}|$ at a location where
the local Unruh temperature is $T_{U}$, the entropy held in escrow against
that temperature is

\begin{equation}
  S_{\mathrm{esc}} = \frac{|U_{\mathrm{grav}}|}{T_{U}}
  \label{eq:6}
\end{equation}

For a test mass m at radius r outside mass M, the binding energy is
$|U_{\mathrm{grav}}|$ = GMm/r in the Newtonian limit, and the proper acceleration
required to hold the test mass static is a = GM/r². The Unruh
temperature at this acceleration is

\begin{equation}
  T_{U} = \frac{\hbar a}{2\pi c k_{B}} = \frac{\hbar GM}{2\pi c k_{B} r^{2}}
  \label{eq:7}
\end{equation}

Substituting (7) into (6):

\begin{equation}
  S_{\mathrm{esc}} = \frac{GMm/r}{\hbar GM/(2\pi c k_{B} r^{2})} = \frac{2\pi c k_{B} r m}{\hbar}
  \label{eq:8}
\end{equation}

This is a real geometric quantity: $S_{\mathrm{esc}}$/$k_{B}$ = 2πr/$\lambda_{C}$ where $\lambda_{C}$ =
$\hbar$/(mc) is the reduced Compton wavelength of the test mass --- that is,
the number of Compton wavelengths that fit into the circumference 2πr.
The dimensionless escrow counts how many Compton circumferences of the
test mass fit around a circle of radius r. This is the quantity the
escrow postulate identifies as the entropy held in reserve by the
gravitational binding.

The same dimensionless quantity has a separate physical meaning that is
worth stating explicitly. The Bekenstein bound states that the entropy
of any system of total energy E confined within a sphere of radius R is
bounded above by

\[ S \le  2πER/(ℏc) \]

For a system of energy E = mc² and confinement radius R = r, the bound
takes the value

\[ 2π(mc²)r/(ℏc) = 2πmcr/ℏ = 2πr/λ_C\qquad(\text{in units of }k_{B}) \]

--- numerically identical to the escrow entropy of a test mass of rest
mass m at orbital radius r given by equation (8). We are explicit that
this is in the first instance a dimensional coincidence: the Bekenstein
bound is a statement about a confined system, and the escrow translation
considers an orbital configuration; reading the equality as 'the escrow
entropy saturates the Bekenstein bound' requires identifying orbital
radius with confinement radius, which is an interpretive step rather
than a settled physical equivalence. With that interpretive
identification made, the equality has a suggestive reading:
gravitational binding holds the test particle's entropy budget at the
same value that the Bekenstein bound would assign to a system of its
rest energy and orbital scale. The framework treats this reading as a
physical motivation, independent of the postulate's dimensional
construction, for considering $S_{\mathrm{esc}}$ rather than some other dimensional
combination as the relevant entropy; the dimensional coincidence itself
is real, and its interpretive content is what the framework proposes.

Because this dimensionless ratio is the candidate organizing variable
for the test-mass leg of the translation, we introduce a symbol for it.
Let

\[ 𝒩_{\mathrm{esc}} \equiv  2πr/λ_C = 2πmcr/ℏ \qquad(8') \]

denote the dimensionless escrow count of a test mass m at radius r.
$\mathcal{N}_{\mathrm{esc}}$ is the number of reduced Compton wavelengths of the test mass that
fit into the circumference of a circle of radius r, and equivalently ---
by the saturation observation above --- the Bekenstein-bound entropy of
the test mass at that radius in units of $k_{B}$. Within the framework,
$\mathcal{N}_{\mathrm{esc}}$ serves as the candidate organizing variable for the test-mass
thermodynamic-gravity dictionary: every other quantity in §II is
expressed in terms of $\mathcal{N}_{\mathrm{esc}}$ and one local kinematic factor ($T_{U}$ or m
c²). The naming is provisional; the dimensionless ratio it labels has
appeared implicitly in semiclassical-gravity calculations for decades
(e.g., in Bekenstein-bound saturation arguments, Euclidean horizon
periodicity, and modular flow generators), but its identification as the
escrow count of a gravitationally bound test mass is, to our knowledge,
first stated in the framework paper [1] and made explicit here. If the
framework is eventually adopted, the dimensionless escrow count is the
most natural quantity to receive a permanent name; we have used
'$\mathcal{N}_{\mathrm{esc}}$' as a working notation rather than attaching a personal label,
in deference to the community's prerogative to settle such naming
through citation rather than authorial declaration.

We can now rewrite the Newtonian-limit time dilation (equation 2) in
escrow language. Multiplying numerator and denominator on the right side
of (2) by $T_{U}$ and recognizing $|U_{\mathrm{grav}}|$/$T_{U}$ = $S_{\mathrm{esc}}$:

\begin{equation}
  \frac{d\tau}{dt} \approx 1 - \frac{GM}{rc^{2}} = 1 - \frac{|U_{\mathrm{grav}}|}{mc^{2}} = 1 - \frac{S_{\mathrm{esc}} T_{U}}{mc^{2}}
  \label{eq:9}
\end{equation}

This is the translation. The fractional slowdown of a clock at radius r
is the product of the local escrow entropy and the local Unruh
temperature, divided by the rest energy of the clock.

Equation (9) is not a new equation. It is equation (2) --- the Newtonian
limit of the Schwarzschild redshift --- with the binding energy
$|U_{\mathrm{grav}}|$ replaced by its escrow decomposition $S_{\mathrm{esc}}$ · $T_{U}$. Its only
content beyond equation (2) is interpretive: it identifies why the clock
slows down at exactly the rate it does. The clock slows because a clock
is an apparatus for converting energy into distinguishable states at a
rate set by its rest mass, and the energy available for
state-distinction at radius r is reduced by the amount $|U_{\mathrm{grav}}|$ that
is held in escrow against the local Unruh temperature.

Numerical verification: at the Earth's equatorial surface (M = 5.972 ×
10²⁴ kg, r = 6.378 × 10⁶ m using WGS84), the right side of equation (9)
evaluates to a fractional clock slowdown of 6.953 × 10⁻¹⁰, which matches
the published Schwarzschild redshift exactly. The corresponding
accumulated clock difference is approximately 22 milliseconds per year,
consistent with measured values.

***C. Why the Tolman shift and the clock slowdown are the same
equation***

The Tolman relation (equation 4) and the redshift relation (equation 1)
are the same equation because of a basic thermodynamic identity. A
thermal bath in equilibrium across a static spacetime is, by definition,
a configuration in which energy flows freely between regions without net
change. For this to be consistent with gravitational redshift --- which
causes photons emitted in a region of low gravitational potential to
lose energy as they climb out --- the local temperature must increase to
compensate. Tolman (1930) derived equation (3) from this requirement.
Equivalently, the Tolman relation is the statement that a thermal
photon's energy redshifts according to (1) and a thermal photon's
temperature is its characteristic energy.

In the escrow vocabulary the identity is immediate. The local Unruh
temperature $T_{U}$(r) is fixed by the proper acceleration a(r) = GM/r². The
escrow held against that temperature is $|U_{\mathrm{grav}}|$/$T_{U}$. The fractional
clock slowdown is $|U_{\mathrm{grav}}|$/(mc²). Eliminating $|U_{\mathrm{grav}}|$ between these
two relations immediately yields

\begin{equation}
  \frac{d\tau}{dt} = 1 - \frac{S_{\mathrm{esc}}\, T_{U}}{mc^{2}} = 1 - \frac{T_{U}\, S_{\mathrm{esc}}}{E_{\mathrm{rest}}}
  \label{eq:10}
\end{equation}

so the clock slowdown is set by the ratio of (local Unruh temperature
times local escrow) to the clock's rest energy. The Tolman relation, in
turn, is the statement that $T_{\mathrm{local}}$ scales as the inverse of the clock
slowdown factor. Both statements reduce to the same algebraic identity:
the local Unruh temperature, the local clock rate, and the local escrow
are three faces of one quantity, namely the local proper acceleration
a(r).

In standard GR this identity is established through the metric: √(−$g_{tt}$)
appears in the redshift, in the Tolman factor, and in the Unruh
temperature through its relationship to surface gravity. The escrow
framework makes the same identity transparent by routing all three
through $S_{\mathrm{esc}}$ and $T_{U}$ explicitly, without requiring the metric as an
intermediate object.


\subsection{What this section does and does not establish}

This section does not derive equation (1) from the escrow postulate.
Equation (1) is a standard result of general relativity and has many
derivations independent of any thermodynamic interpretation. What this
section establishes is a dictionary: the standard expression for
gravitational time dilation can be rewritten in escrow language as
equation (10), and the rewriting makes manifest the relationship between
time dilation and the Tolman temperature shift that is otherwise an
apparent coincidence requiring separate derivation. We note explicitly
that the escrow translation in equation (9) operates at the level of the
weak-field expansion (1 − GM/rc²), not the exact Schwarzschild factor
√(1 − 2GM/rc²). It is important to distinguish two separate statements:
the Tolman--Ehrenfest relation $T_{\mathrm{local}}$ · √(−$g_{tt}$) = $T_{\infty}$, equation (3),
is exact for any static spacetime (it follows from the existence of a
timelike Killing vector and from Tolman--Ehrenfest photon-redshift
equilibrium); the weak-field approximation enters only when one inserts
the linearized Schwarzschild metric factor as √(1 − 2GM/rc²) ≈ 1 −
GM/rc². The escrow translation of the Schwarzschild redshift in equation
(9) is therefore weak-field, but the underlying Tolman relation is not.
The exact Schwarzschild form can be obtained by inserting the full
metric factor into the Tolman relation rather than by extending the
escrow translation itself. The framework's structural commitment at
this level is to the weak-field reading of equation (9), with the
strong-field exact form available through the (exact) Tolman relation.

The escrow framework's contribution at this level is interpretive, not
predictive. No measurement of time dilation distinguishes the escrow
reading from the standard GR reading; the equations are identical. The
framework's value, here, is a conceptual re-reading: the postulate
$S_{\mathrm{esc}}$ = $|U_{\mathrm{grav}}|$/$T_{U}$ provides a single object --- in the
partly-notational sense developed in §V.G--H --- through which time
dilation, Tolman temperature, Unruh temperature, and gravitational
binding energy are all expressed, replacing four separately-stated
relations with one.

This conceptual re-reading is what we will exploit in Section III, where
we apply the same dictionary to the Jacobson derivation of Einstein's
equations and the Bekenstein--Hawking entropy.


\subsection{A note on what the framework does not commit to}

We do not invoke a "vacuum computational bandwidth" or "information
processing rate" reading of equation (10). Such readings appear in some
popularizations of entropic gravity, and they are not part of the escrow
postulate as stated in [1]. The framework's commitment is to the
thermodynamic identity (6) and its consequences; it does not commit to a
digital or computational substrate for the vacuum. The local rate of
escrow production at radius r is a thermal quantity, set by $T_{U}$(r), and
is correctly described in the vocabulary of equilibrium thermodynamics.
No additional ontological structure is required to make equation (10)
hold.

Similarly, we do not introduce a separate "escrow stress-energy
tensor" T\textsubscript{μν}\textsuperscript{escrow} to be added to the matter stress-energy on the
right side of Einstein's equation. The reason is structural:
Einstein's equations are an algebraic identity between the matter
stress-energy tensor T\textsubscript{μν}\textsuperscript{matter} and the geometric Einstein tensor
G\textsubscript{μν}, and the Bianchi identities (∇\textsuperscript{μ} G\textsubscript{μν} = 0) require ∇\textsuperscript{μ}
T\textsubscript{μν}\textsuperscript{matter} = 0 --- the matter content already saturates the structure
of conserved sources. Gravitational binding energy is not a localizable
tensorial quantity in general relativity; it is encoded non-locally in
the geometry itself (the difference, for an isolated system, between the
ADM mass at infinity and the integrated rest mass on a spatial slice),
and it cannot be added to T\textsubscript{μν}\textsuperscript{matter} as an independent source. The
escrow framework's $S_{\mathrm{esc}}$ is a thermodynamic interpretation of a
configuration's gravitational binding content; it does not introduce a
new tensor and does not augment Einstein's equation with any new
stress-energy contribution. Attempting to do so would either violate the
Bianchi identities or double-count: the geometry on the left of
Einstein's equation already responds to the full gravitational
configuration, and a separate escrow source would constitute a second
response to the same physical content. This point will be developed
further in Section III.


\section{The Jacobson Translation}


\subsection{The standard statement}

In 1995, Jacobson showed [3] that Einstein's field equations follow
from demanding that the thermodynamic relation $\delta Q$ = T·dS hold across
every local Rindler horizon, where T is the Unruh temperature of an
observer accelerating just outside the horizon and dS is the area-law
entropy of the horizon. The derivation is now standard and we
recapitulate only the load-bearing steps. Consider a point p in
spacetime and a small local Rindler horizon through p, defined by the
future causal horizon of an accelerated observer with proper
acceleration a passing through p. The Unruh temperature associated with
this horizon is

\begin{equation}
  T_{U} = \frac{\hbar a}{2\pi c k_{B}}
  \label{eq:11}
\end{equation}

The entropy associated with the horizon is its area in Planck units,

\begin{equation}
  S_{\mathrm{horizon}} = \frac{A}{4\ell_{P}^{2}}\, k_{B}
  \label{eq:12}
\end{equation}

where $\ell_{P}$ = ($\hbar$G/c³)\textsuperscript{1/2} is the Planck length. A small area element of
the horizon expands or contracts as matter falls into the wedge behind
it, and the rate of expansion is governed by the Raychaudhuri equation
applied to the null geodesics generating the horizon.

Jacobson's central insight was that demanding

\begin{equation}
  \delta Q = T_{U}\, dS_{\mathrm{horizon}}
  \label{eq:13}
\end{equation}

hold at every such point p, for every choice of local Rindler horizon,
forces the Einstein field equations on the spacetime:

\begin{equation}
  R_{\mu\nu} - \tfrac{1}{2} g_{\mu\nu} R + \Lambda g_{\mu\nu} = \frac{8\pi G}{c^{4}}\, T_{\mu\nu}
  \label{eq:14}
\end{equation}

where T\textsubscript{μν} is the matter stress-energy tensor whose flux $\delta Q$ across the
horizon enters the left side of (13). The Raychaudhuri equation supplies
the geometric content; the Clausius relation $\delta Q$ = T·dS supplies the
thermodynamic content; the combination is Einstein's equation. This
derivation is not specific to Schwarzschild or to any particular
solution --- it holds locally everywhere, for every observer.

Jacobson's result establishes that Einstein's equation has the
standing of a thermodynamic equation of state of the vacuum, analogous
to the ideal-gas law PV = NkT emerging from molecular kinetics rather
than being postulated fundamentally. The metric, in this reading, is
what you get when the vacuum's entropy obeys the second law locally
everywhere; it is not a separate object whose dynamics need to be
derived from a Lagrangian.


\subsection{The escrow translation}

The static escrow postulate identifies the entropy held against the
local Unruh temperature as the gravitational binding energy divided by
that temperature: $S_{\mathrm{esc}}$ = $|U_{\mathrm{grav}}|$/$T_{U}$ (equation 6). Jacobson's
derivation uses T = $T_{U}$ and S = $S_{\mathrm{horizon}}$ (the horizon area entropy).
The question is: what is the relationship between $S_{\mathrm{horizon}}$ as it
appears in Jacobson's derivation and $S_{\mathrm{esc}}$ as it appears in the escrow
postulate?

The answer is that they are the same quantity, evaluated at different
scales. Jacobson's $S_{\mathrm{horizon}}$ is the entropy associated with a
particular local Rindler horizon --- a coarse-grained quantity defined
at the level of the horizon's area. The escrow entropy $S_{\mathrm{esc}}$ is the
entropy held against the same Unruh temperature by whatever mass-energy
is in the wedge behind that horizon. The first law of horizon
thermodynamics, applied to a small flux of mass-energy crossing the
horizon, equates these two: a flux $\delta Q$ of energy at temperature $T_{U}$
across the horizon produces a change $\delta S$ in the horizon entropy, and that
$\delta S$ is exactly $\delta Q$/$T_{U}$ --- the change in escrow entropy associated with the flux. (The naive substitution $\delta Q$/$T_{U}$ = δ$|U_{\mathrm{grav}}|$/$T_{U}$ would be off by a factor of two under $|U_{\mathrm{grav}}| = Mc^{2}/2$; the escrow differential $dS_{\mathrm{esc}} = d(|U_{\mathrm{grav}}|/T_{U})$ reproduces $\delta Q$/$T_{U}$ correctly because $T_{U}\propto 1/M$ makes the temperature-variation term contribute equally, whereas the one-term reading δ$|U_{\mathrm{grav}}|$/$T_{U}$ alone omits it.) We note that the first law relates the two
entropies at the differential level (δ$S_{\mathrm{horizon}}$ = δ$S_{\mathrm{esc}}$); the finite
identification $S_{\mathrm{horizon}}$ = $S_{\mathrm{esc}}$ for a Schwarzschild black hole of mass
M follows from this differential relation together with integrability
along the formation path, and is confirmed directly by the algebraic
calculation in §III.C (equations 17--18). The finite-level confirmation
does not depend solely on the differential argument.

In escrow language, Jacobson's identity (13), with the substitution
$S_{\mathrm{horizon}}$ → $S_{\mathrm{esc}}$ justified at the differential level by the first law
and confirmed at the finite level by §III.C below, reads

\begin{equation}
  \delta Q = T_{U}\, dS_{\mathrm{esc}}
  \label{eq:15}
\end{equation}

which is the differential version of the static escrow postulate (6).
The static postulate says $S_{\mathrm{esc}}$ = $|U_{\mathrm{grav}}|$/$T_{U}$; the differential form
says that the change in escrow entropy across a horizon equals the heat
flux across that horizon divided by the local Unruh temperature. These
are the same statement. Equation (15) is therefore not a new equation.
It is Jacobson's (13) with $S_{\mathrm{horizon}}$ replaced by $S_{\mathrm{esc}}$, justified by
the observation that the horizon entropy is exactly the entropy held
against the horizon's Unruh temperature --- which is the escrow
definition.

The substitution buys nothing mathematically. It buys an interpretation:
what the horizon is doing, thermodynamically, is bookkeeping the
gravitational binding energy of the matter behind it. The horizon area
grows when matter falls in because the matter brought additional binding
energy to the wedge, and that additional binding energy must be held in
escrow against the horizon's Unruh temperature. The escrow framework
supplies the answer to the question "what is the horizon's entropy the
entropy of?" --- it is the entropy held in escrow against the local
Unruh temperature by the gravitationally bound mass-energy in the wedge.


\subsection{Why the substitution is legal}

The identification $S_{\mathrm{horizon}}$ = $S_{\mathrm{esc}}$ requires that two independently
defined quantities coincide. The first is the horizon area entropy
A/(4$\ell_{P}$²) · $k_{B}$, which is a property of the horizon geometry alone. The
second is the escrow entropy $|U_{\mathrm{grav}}|$/$T_{U}$, which is a property of the
matter content and the local Unruh temperature. We now show that these
are equal.

Consider a Schwarzschild black hole of mass M. The Schwarzschild radius
and surface gravity at the horizon are

\[ r_s = 2GM/c², κ = c⁴/(4GM) \]

The Hawking temperature --- the Tolman-redshifted (asymptotic) Unruh
temperature of an observer hovering near the horizon, finite at infinity
even though the local Unruh temperature of a static observer diverges as
r → $r_{s}$, with

\[ T_H = lim_{r\to r_s} T_{\mathrm{Unruh}}(r) \cdot  \sqrt{-g_{\mathrm{tt}}(r)} \]

--- is

\begin{equation}
  T_{U} = \frac{\hbar \kappa}{2\pi c k_{B}} = \frac{\hbar c^{3}}{8\pi G M k_{B}}
  \label{eq:16}
\end{equation}

The self-binding energy of the configuration at the horizon scale is
$|U_{\mathrm{grav}}|$ = GM²/$r_{s}$ = Mc²/2 (a direct consequence of the Schwarzschild
radius definition). This value also coincides with the Schwarzschild (Q
= J = 0) limit of the Smarr identification T\textsubscript{H} $S_{\mathrm{BH}}$ = (M − ΦQ − 2Ω\textsubscript{H}
J)/2 used in §III.D below, which gives the Mc²/2 choice a thermodynamic
motivation --- it is the gravitational sector of the Smarr
decomposition, with the electrostatic and rotational sectors vanishing
for an uncharged non-rotating horizon --- alongside the algebraic
GM²/$r_{s}$ motivation. This choice of $|U_{\mathrm{grav}}|$ is one of the four
regime-specific identifications enumerated in §V.G; it is not derived
from a single covariant binding-energy construction (such a construction
is identified there as an open theoretical task) but is adopted because
it produces the algebraic match with $S_{\mathrm{BH}}$ that the present subsection
demonstrates. The escrow entropy is therefore

$S_{\mathrm{esc}}$/$k_{B}$ = $|U_{\mathrm{grav}}|$/($k_{B}$ $T_{U}$) = (Mc²/2) · (8πGM$k_{B}$)/($\hbar$c³ $k_{B}$) =
\begin{equation}
  \frac{4\pi G M^{2}}{\hbar c} \equiv \mathcal{N}_{\mathrm{esc}}^{\mathrm{Schw}}
  \label{eq:17}
\end{equation}

where we have introduced the horizon analog of the dimensionless escrow
count --- $\mathcal{N}_{\mathrm{esc}}^{\mathrm{Schw}}$ ≡ 4πGM²/($\hbar$c) = 4π(M/$m_{P}$)², with $m_{P}$ = √($\hbar$c/G) the
Planck mass --- in analogy with the test-mass count $\mathcal{N}_{\mathrm{esc}}$ ≡ 2πr/$\lambda_{C}$ of
equation (8') in §II.B. The horizon count $\mathcal{N}_{\mathrm{esc}}^{\mathrm{Schw}}$ measures the
gravitationally bound mass-energy of a Schwarzschild black hole in units
of the Planck-scale escrow quantum. We are explicit that $\mathcal{N}_{\mathrm{esc}}$ and
$\mathcal{N}_{\mathrm{esc}}^{\mathrm{Schw}}$ are not the same physical quantity in the literal sense: the
test-mass $\mathcal{N}_{\mathrm{esc}}$ has the geometric interpretation 'reduced Compton
wavelengths fitting into the orbital circumference' (a mode count in
the sense of counting wavelengths around a closed curve), while
$\mathcal{N}_{\mathrm{esc}}^{\mathrm{Schw}}$ is the Bekenstein--Hawking entropy in dimensionless form
(scaling as M² rather than as the linear mr of the test-mass case). They
belong to the same family of dimensionless escrow constructions --- each
arising from the postulate S = |U|/T applied in a different regime ---
but the family is bound by the common postulate, not by a single
geometric mode-count formula. The shared notation marks the structural
relationship; the regime label ('\textsuperscript{Schw}') marks the difference in
physical content. The Bekenstein--Hawking entropy of the same horizon is

\begin{equation}
  \frac{S_{\mathrm{BH}}}{k_{B}} = \frac{A}{4\ell_{P}^{2}} = \frac{4\pi r_{s}^{2} c^{3}}{4\hbar G} = \frac{4\pi G M^{2}}{\hbar c} = \mathcal{N}_{\mathrm{esc}}^{\mathrm{Schw}}
  \label{eq:18}
\end{equation}

Equations (17) and (18) are identical. The escrow entropy of a
Schwarzschild black hole, computed from the self-binding energy and the
Hawking temperature via the static escrow postulate, equals the
Bekenstein--Hawking entropy exactly. There is no approximation and no
fudge factor; the geometric 4π in the horizon area is recovered exactly
by the geometric 4π in the Hawking temperature formula. This is the
content of the static escrow postulate's recovery of
Bekenstein--Hawking, originally reported in the framework paper [1],
restated here in the Jacobson-translation context.

The Bekenstein--Hawking formula by itself supplies the relation between
area and entropy without supplying an interpretation of what the entropy
is the entropy of. The escrow framework supplies that interpretation:
the horizon entropy is the entropy held in escrow against the local
Unruh temperature by the gravitationally bound mass-energy inside the
horizon. This is a structural interpretive claim, equivalent in
mathematical content to Bekenstein--Hawking but distinct in physical
reading. The horizon's area-as-entropy is not a coincidence of
dimensional analysis nor an axiom requiring separate justification ---
it is the escrow of the matter trapped behind the horizon, evaluated at
the horizon's own Unruh temperature.


\subsection{Extension to Reissner--Nordström and Kerr horizons}

A natural question for the escrow postulate, raised in adversarial
review of this paper, is whether the Schwarzschild recovery in equations
(17)--(18) is the only horizon entropy the postulate gets right, or
whether the same form $S_{\mathrm{esc}}$ = $|U_{\mathrm{grav}}|$/T\textsubscript{H} continues to recover
Bekenstein--Hawking entropies for charged (Reissner--Nordström) and
rotating (Kerr) black holes. We report the result of that test, with
explicit attention to what is and is not established by it.

For a Reissner--Nordström black hole of mass M and charge Q (in
geometrized units G = c = k\textsubscript{e} = 1, with k\textsubscript{e} = 1/(4πε₀) the Coulomb
constant), the outer and inner horizons, area, and Hawking temperature
are

\[ r_\pm  = M \pm  \sqrt{M² - Q²}, A = 4πr_+², T_H = (r_+ - r_-)/(4πr_+²) \]

The Smarr formula relates these to the mass:

\[ M = 2T_H S_{\mathrm{BH}} + Φ Q \]

where Φ = Q/r\textsubscript{+} is the electrostatic potential at the horizon and $S_{\mathrm{BH}}$
= A/4 is the Bekenstein--Hawking entropy. Rearranged,

\[ T_H S_{\mathrm{BH}} = (M - Φ Q)/2 = (M - Q²/r_+)/2 \]

For the escrow postulate $S_{\mathrm{esc}}$ = $|U_{\mathrm{grav}}|$/T\textsubscript{H} to recover $S_{\mathrm{BH}}$ exactly,
the gravitational binding energy must satisfy

$|U_{\mathrm{grav}}|$\textsuperscript{RN} = (M − Q²/r\textsubscript{+})/2 (geometrized); (Mc²/2) − k\textsubscript{e} Q²/(2r\textsubscript{+})
(SI)

i.e. the Schwarzschild binding minus the electrostatic energy stored at
the horizon.

For Kerr (mass M, angular momentum J = Ma), the horizon radius, area,
and Hawking temperature are

r\textsubscript{+} = M + √(M² − a²), A = 4π(r\textsubscript{+}² + a²), T\textsubscript{H} = (r\textsubscript{+} −
r\textsubscript{−})/(4π(r\textsubscript{+}² + a²))

and the Smarr formula reads

\[ M = 2T_H S_{\mathrm{BH}} + 2Ω_H J \]

where Ω\textsubscript{H} = a/(r\textsubscript{+}² + a²) is the angular velocity of the horizon.
Rearranged,

\[ T_H S_{\mathrm{BH}} = (M - 2Ω_H J)/2 \]

The required Kerr binding energy is therefore

\[ |U_{\mathrm{grav}}|^{\mathrm{Kerr}} = (Mc²/2) - Ω_H J \]

the Schwarzschild binding minus the rotational work term.

Both extensions are exact: with these natural identifications of
$|U_{\mathrm{grav}}|$, the escrow postulate recovers $S_{\mathrm{BH}}$ for RN and Kerr to all
displayed digits across the parameter range, including near extremality
(where T\textsubscript{H} → 0 and $S_{\mathrm{esc}}$ takes the 0/0 form whose limit is the finite
extremal entropy). The structural pattern is consistent across all three
cases: $|U_{\mathrm{grav}}|$ equals the Schwarzschild gravitational binding minus
the matter-field chemical-potential terms (Φ Q for charge, Ω J for
rotation) that appear in the first law of black hole mechanics.

We are explicit about the interpretive status of these extensions. They
are structural-consistency checks rather than independent predictions.
The Smarr formula, which is itself derivable from the area law and the
first law, already specifies T\textsubscript{H} $S_{\mathrm{BH}}$ as (M − ΦQ − 2Ω\textsubscript{H} J)/2; the escrow
postulate's apparent "recovery" of $S_{\mathrm{BH}}$ for RN and Kerr is the
algebraic statement that $|U_{\mathrm{grav}}|$ is identified with this same
quantity. The framework does not independently specify $|U_{\mathrm{grav}}|$ for
charged or rotating configurations and then derive $S_{\mathrm{BH}}$; it adopts the
chemical-potential decomposition that the first law already supplies,
and reads off the consistency. What the consistency check does establish
is that the escrow postulate's form $S_{\mathrm{esc}}$ = $|U_{\mathrm{grav}}|$/T\textsubscript{H} survives the
addition of matter fields and rotation without modification, which is a
non-trivial structural fact: a postulate that produced $S_{\mathrm{BH}}$ only for
Schwarzschild and failed for RN or Kerr would be ruled out by these
calculations. The escrow framework is therefore consistent with the full
first law of black hole mechanics, not just with the Schwarzschild
special case.


\subsection{What this section does and does not establish}

This section does not derive Einstein's field equations from the escrow
postulate. Jacobson did that in 1995 [3], using $\delta Q$ = T·dS as input.
What this section establishes is that the S appearing in Jacobson's
input has the interpretation S = $S_{\mathrm{esc}}$, with the consequence that the
derived Einstein equation is the equation of state of the vacuum's
escrow thermodynamics.

Two consequences follow. First, the escrow framework does not require
any modification of Einstein's equations; it provides an interpretation
of why those equations have the structure they do. The cosmological
constant Λ appears on the left side of (14) without special handling
because Jacobson's derivation produces it as an integration constant in
the same way it produces the rest of the geometry. Second, attempts to
add a separate "escrow stress-energy tensor" to the right side of
Einstein's equation are explicitly ruled out by this derivation: the
entire content of $\delta Q$ on the right side of Jacobson's identity is
already the matter stress-energy flux T\textsubscript{μνk}\textsuperscript{μk}\textsuperscript{νdA} dλ (where k\textsuperscript{μ} is
the null generator of the horizon and λ is the affine parameter). The
escrow contribution lives entirely in the interpretation of the S on the
left side, not in any additional source on the right.

The framework's value at this level is the same as it was for time
dilation: a conceptual re-reading. Jacobson's derivation, the
Bekenstein--Hawking formula, the Tolman temperature law, the Unruh
effect, and the static escrow postulate are five separately-derived
results that the escrow framework expresses through a common
thermodynamic object $S_{\mathrm{esc}}$. None of these results is modified; their
interrelation is made transparent.


\subsection{A note on what the framework does not commit to}

We follow Jacobson in treating equation (14) as "an" equation of state
rather than asserting that spacetime is literally a thermodynamic system
in some fundamental sense. The thermodynamic analogy is structural ---
the same Clausius relation, the same area-as-entropy identification, the
same local equilibrium condition --- not ontological. We do not claim
that spacetime is built out of more fundamental thermodynamic degrees of
freedom in the way that an ideal gas is built out of molecules. Such a
claim would require identifying the microscopic constituents whose
statistical mechanics produces the area law; this is the project of
quantum gravity and is not addressed by the escrow framework.

Similarly, we do not claim that the entropy in equation (15) is
observer-independent. The Unruh temperature depends on the choice of
accelerated observer; the horizon depends on the choice of local Rindler
frame; the escrow entropy inherits both dependences. What is
observer-independent is the relation $\delta Q$ = $T_{U}$ · d$S_{\mathrm{esc}}$, which holds for
every observer in their own frame. This observer-dependence is a feature
of the thermodynamic interpretation, not a bug: it is what distinguishes
Jacobson's reading of GR from a global-entropy interpretation in which
a unique S is assigned to a spacetime. Different observers see different
escrow entropies because they have different local Unruh temperatures;
the field equation (14) is what enforces consistency across observers.

The escrow framework's identification of $S_{\mathrm{horizon}}$ with $S_{\mathrm{esc}}$ therefore
does not add a substantive claim beyond what Jacobson's derivation
already commits to. It supplies a name and an interpretation for the
entropy that Jacobson uses as input.


\section{The Modular Hamiltonian Connection}


\subsection{The standard statement}

Dimensional bookkeeping for this section: throughout §IV we follow the
standard entanglement-entropy convention in which the modular
Hamiltonian K is dimensionless (it appears as the exponent in the
reduced density matrix ρ = exp(−K)/Z, so dimensionlessness is required
by the exponential's argument). Entanglement entropy S\textsubscript{EE} in this
convention is also dimensionless (a count, measured in nats or
equivalently in units of $k_{B}$). Where this section writes $T_{U}$ · $\Delta S_{\mathrm{EE}}$ or
$T_{U}$ · $\Delta\langle K\rangle$, the implicit factor of $k_{B}$ converting dimensionless entropy
to thermodynamic entropy is understood: the energy associated with a
modular-Hamiltonian-content change $\Delta\langle K\rangle$ at Unruh temperature $T_{U}$ is $k_{B}$
$T_{U}$ · $\Delta\langle K\rangle$, exactly as in §IV.A's calculation ΔK · $k_{B}$ · $T_{U}$(d1) = mc².
Sections II and III used the alternative convention in which S is
written in J/K from the start (with $k_{B}$ explicit in each formula); the
two conventions differ only by where the $k_{B}$ is placed, and we use
whichever is cleaner in each setting. Equations (19) and (20) below are
written in the entanglement-entropy convention; explicit $k_{B}$ factors are
inserted where the equation crosses into thermodynamic-entropy
quantities.

Jacobson's derivation [3] uses the horizon area entropy as the input
S in $\delta Q$ = T·dS, with the area law itself taken as a given. A subsequent
body of work, of which the most consequential single result is due to
Faulkner, Guica, Hartman, Myers, and Van Raamsdonk (FGHMV, [5,6]),
establishes a related but distinct identification: in asymptotically AdS
spacetimes, demanding the entanglement-entropy first law $\delta S$\textsubscript{EE} = δ$\langle K\rangle$
for vacuum-state perturbations across spherical entangling surfaces
forces the linearized Einstein equations in the bulk. Here S\textsubscript{EE} is the
entanglement entropy of a region of the boundary CFT, K is the modular
Hamiltonian of that region in the vacuum state (the operator whose
exponential gives the reduced density matrix on the region), and δ$\langle K\rangle$ is
the expectation value of K in the perturbed state minus its value in the
vacuum.

The FGHMV result is structurally similar to Jacobson's but
operationally distinct in two respects. First, it uses entanglement
entropy rather than horizon area as the input quantity, with the
Ryu--Takayanagi formula [10] supplying the bulk-area dual of the
boundary entanglement. Second, it produces linearized Einstein equations
rather than the full nonlinear theory, with the linearization being
around the AdS vacuum. Subsequent work has extended the result to
broader settings, including non-spherical entangling regions and
curved-background perturbations, with the unifying structural statement
being that the gravitational equations of motion are equivalent to the
requirement that relative entropy (a closely related and equally
fundamental information-theoretic quantity) be non-negative for
vacuum-state perturbations.

The vacuum modular Hamiltonian for a Rindler wedge has an exact form,
originally derived by Bisognano and Wichmann [11]:

\[ K = (2π/(ℏc)) \int  x{\perp} T_{00}(x) dV \]

where x⊥ is the perpendicular distance from the cut and T\textsubscript{00} is the
energy density. This is the same 2π that appears in the Unruh
temperature --- the modular flow generated by K is exactly the boost
flow that defines the Rindler wedge, and the period 2π in boost
parameter corresponds to the inverse Unruh temperature. The
Bisognano--Wichmann formula gives a sharp prediction for the change in
modular Hamiltonian content induced by inserting a localized energy E =
mc² at perpendicular distance d1 from the cut:

ΔK = 2π·m·d1 (natural units; $\hbar$ = c = $k_{B}$ = 1)

or equivalently, restored to SI units,

\[ ΔK = 2π\cdot m\cdot c\cdot d1/ℏ = 𝒩_{\mathrm{esc}}(d1) \]

(a dimensionless quantity, as required of a modular Hamiltonian
content), where $\mathcal{N}_{\mathrm{esc}}$(d1) is the escrow count of equation (8')
evaluated at r = d1. The associated energy at the Rindler-observer Unruh
temperature is

\[ T_U(d1) = ℏc/(2π\cdot d1\cdot k_B), ΔK\cdot k_B\cdot T_U(d1) = mc² \]

--- the rest energy of the perturbation. This is the prediction that the
lattice paper [2] tests numerically in free scalar field theory on
regulated 1+1D and 3+1D lattices.


\subsection{The escrow translation}

The escrow postulate identifies $S_{\mathrm{esc}}$ = $|U_{\mathrm{grav}}|$/$T_{U}$ as the entropy
held against the local Unruh temperature by gravitationally bound
mass-energy. The FGHMV result uses S\textsubscript{EE} and the modular Hamiltonian K,
neither of which is identical to $S_{\mathrm{esc}}$ as defined. The question is: what
is the relationship between the escrow entropy and the modular
Hamiltonian content?

The relationship is mediated by the local first law. For a vacuum-state
perturbation localized in a Rindler wedge, the entanglement-entropy
first law $\delta S$\textsubscript{EE} = δ$\langle K\rangle$ says that the change in entanglement entropy of
the wedge (in dimensionless units; see the unit-convention statement at
the top of §IV.A) equals the change in expectation value of the wedge's
modular Hamiltonian. The product $k_{B}$ · $T_{U}$ · $\delta S$\textsubscript{EE} has units of energy
and represents the heat content associated with the entanglement change.
Identifying this with the gravitational binding energy of the
perturbation gives

\begin{equation}
  |U_{\mathrm{grav}}| = k_{B}\, T_{U}\, \Delta S_{\mathrm{EE}} = k_{B}\, T_{U}\, \Delta\langle K\rangle
  \label{eq:19}
\end{equation}

which, by the escrow postulate (6) $S_{\mathrm{esc}}$ = $|U_{\mathrm{grav}}|$/$T_{U}$, means

\begin{equation}
  S_{\mathrm{esc}} = k_{B}\, \Delta S_{\mathrm{EE}} = k_{B}\, \Delta\langle K\rangle
  \label{eq:20}
\end{equation}

with the explicit $k_{B}$ factor converting the dimensionless
entanglement-entropy quantities of the entanglement-gravity literature
into the thermodynamic-entropy units in which $S_{\mathrm{esc}}$ was defined in §II.
The dimensionless content of equation (20) is $S_{\mathrm{esc}}$/$k_{B}$ = $\Delta S_{\mathrm{EE}}$ = $\Delta\langle K\rangle$,
which is what the lattice computations of [2] compute directly.

We are explicit about the regime-bridging assumptions that this
identification carries. The FGHMV result [5,6] is formulated for
asymptotically AdS spacetimes with spherical entangling surfaces on the
boundary CFT, with the relevant entropy being the entanglement entropy
of a boundary subregion and the relevant modular Hamiltonian being the
operator generating modular flow for the vacuum reduced density matrix
of that subregion. The escrow framework's identification in (19)--(20)
applies the same first-law structure to a localized perturbation in a
flat-space Rindler wedge, where $T_{U}$ is the Unruh temperature of an
accelerated observer. The bridge between these two settings carries
three distinct technical assumptions, each well-established in the
entanglement-entropy literature [5,6,11,12,13] but each worth stating
individually: (a) the modular-Hamiltonian content of a Rindler wedge in
Minkowski vacuum is given by the Bisognano--Wichmann formula K =
(2π/($\hbar$c)) ∫ x⊥ T\textsubscript{00} dV, where x⊥ is perpendicular distance from the cut;
(b) the entanglement-first-law statement $\delta S$\textsubscript{EE} = δ$\langle K\rangle$ carries over from
the boundary CFT setting to a flat-space Rindler wedge in the
linear-response regime, with no curvature corrections to the first law
at this order; (c) the gravitational binding energy of the localized
perturbation, identified with the source energy mc², plays the same role
on the matter side that boundary-CFT energy density plays in the AdS
setting. Each of (a), (b), (c) is standard in the literature but each is
an additional assumption beyond the bare FGHMV statement. We treat all
three as assumptions rather than as derived statements, and the
framework's modular-Hamiltonian leg therefore inherits whatever caveats
attend each of these standard identifications in the broader
entanglement-gravity literature.

The escrow entropy of a localized mass perturbation, under this
identification, equals both the change in entanglement entropy of the
wedge containing it and the change in expectation value of the wedge's
modular Hamiltonian. Equation (20) is the strongest interpretive content
of the escrow framework: the escrow entropy is the modular-Hamiltonian
content of a localized gravitational perturbation, with the FGHMV result
then identifying the consistency requirement on this content with
linearized Einstein equations.

If equation (20) held exactly, the modular Hamiltonian content of a
localized mass m at perpendicular distance d1 from the cut would equal
the escrow count $\mathcal{N}_{\mathrm{esc}}$(d1) = 2π·m·c·d1/$\hbar$ (the dimensionless content of
equation 8' with r = d1). The Bisognano--Wichmann formula gives exactly
this. Equation (20) therefore predicts that lattice computations of ΔK
in regulated free scalar field theory should yield the BW linear
asymptote with prefactor unity, i.e. ΔK = $\mathcal{N}_{\mathrm{esc}}$(d1) at small d1 in the
half-space limit. This is the prediction the lattice paper [2] tests
numerically.


\subsection{What the lattice tests}

The lattice paper [2] computes ΔK directly on regulated free scalar
lattices in 1+1D (chain, N up to 3000) and 3+1D (cubic, N up to 14)
using the Williamson-decomposition method, with two protocols: a
two-mass protocol with masses at perpendicular distances d1 and d2 from
the cut, and a single-mass-in-A protocol with a single mass at d1 (the
latter being the cleaner test of the BW prediction). The numerical
results, summarized in [2], are:

In 1+1D at small d1 (d1 ∈ [2, 6] at m = 1, N = 1000), the local
power-law exponent of ΔK vs d1 is +1.02 ± 0.05, indistinguishable from
the BW linear asymptote α = 1. The prefactor is approximately
0.030--0.037, corresponding to a suppression of approximately 1/30 of
the literal BW value. At larger d1 the local exponent decreases smoothly
toward zero, reflecting geometric corrections to the BW half-space
assumption (Dirichlet boundary conditions, finite subsystem-A depth).
The peak prefactor is independently reproducible from the lattice
paper's published methodology, with both V1 (relative-entropy
positivity) and V2 (modal-versus-original-basis cross-check) sanity
checks passing at every d1 in the sweep; the analytic origin of the
prefactor remains the open theoretical problem identified in this
section.

In 3+1D at all d1 the lattice can resolve (N ≤ 14, d1 ≤ 6), ΔK(d1) is
peaked at d1 ≈ 2 and decays as $d_1^{-2\text{ to }-3}$ for d1 ≥ 3. The BW linear
asymptote is not recovered at any d1 the lattice can resolve. The peak
structure reproduces across N = 12 and N = 14, indicating that the
result is not a finite-size artifact at the lattice sizes used.

Taken at face value, these results constrain equation (20) in two ways.
In 1+1D, the equation holds approximately within a small-d1 window with
a prefactor suppression of 1/30; the prefactor has no
presently-calculated origin, and finding one is the principal open
theoretical question. In 3+1D, the equation does not hold within the d1
range the lattice can resolve; whether it can be recovered at much
larger d1 inaccessible to the present lattice (corresponding to a slower
rate of approach to the continuum BW asymptote in higher dimensions) is
an open empirical question.


\subsection{What this section does and does not establish}

This section does not establish equation (20) as a derived identity. It
identifies (20) as the natural escrow reading of the FGHMV result, with
[2] supplying empirical content on whether the identification holds
when both sides are computed independently. The 1+1D partial recovery
and the 3+1D non-recovery within accessible d1 are the present state of
evidence; neither rules out (20) fundamentally, and neither establishes
it without qualification.

Three distinct theoretical questions are sharpened by this section.
First, what is the calculable origin of the 1/30 prefactor in 1+1D? A
candidate explanation flagged in [2] is that the lattice-regularized
free scalar at m = 1 has correlation length ξ ∼ 1/m ∼ 1 lattice spacing,
making the localized mass a UV-scale perturbation by the lattice's own
measure, while the BW formula is an infrared / long-wavelength
prediction. The 1/30 suppression may reflect a UV/IR mismatch
correctable by a first-principles calculation of the lattice-regulated
modular Hamiltonian. The Dirichlet boundary conditions used for the
lattice differ from the free BCs assumed in the standard BW derivation,
and the boundary-correction factor is another candidate contribution.
Resolving the prefactor's origin from first principles is the single
most consequential follow-up calculation the framework presently faces.

Second, does the 3+1D non-recovery reflect a genuinely shorter
characteristic length scale of the 3+1D modular Hamiltonian content, or
a slower rate of approach to the continuum BW limit due to the smaller
conformal-symmetry group of the 3+1D massless free scalar? The 1+1D
massless scalar is a CFT with infinite-dimensional Virasoro symmetry;
the 3+1D massless scalar has only the finite-dimensional conformal
group, and the rate at which the lattice-regularized modular content
approaches the continuum BW asymptote may differ between them for this
reason. The 3+1D test on N ≳ 20 lattices (computationally feasible on
dedicated GPU hardware but not on the present authors' CPU resources)
would distinguish these two interpretations.

Third, what is the dimensional-dependence theorem, if any, connecting
1+1D BW partial recovery to the 3+1D structure observed in [2]? A 2+1D
test on the same protocol would be the natural intermediate case and
would determine whether the dimensional dependence is monotone. The
present authors leave this as a future calculation.


\subsection{A note on what the framework does not commit to}

We do not claim that the lattice paper [2]'s partial recovery in 1+1D
establishes equation (20) as a derived identity. The 1/30 prefactor is
an unexplained suppression, not a confirmed match; the framework's
interpretive content at this level depends on that prefactor having a
calculable origin, and the calculation has not been performed. The
framework is at risk of being read as having empirical confirmation it
does not yet have, and we explicitly disclaim that reading: [2] is
consistent with the escrow framework's modular-Hamiltonian
identification at the level of qualitative structure (positive sign,
linear scaling at small d1, the right dimensionless quantity $\mathcal{N}_{\mathrm{esc}}$), but
it is not yet quantitatively confirmed at the level of the absolute
prefactor.

We do not claim that the 3+1D non-recovery rules out the escrow
framework. As stated above, the rate-of-approach question is unresolved
by the present lattice sizes, and a finding that 3+1D BW is recovered at
much larger d1 inaccessible to the present lattice would resolve the
non-recovery as a finite-size effect rather than a fundamental
obstruction. The framework's commitment is to the postulate (6) and the
interpretive content of equations (10), (15), (17), and (20); empirical
confirmation of these at the lattice level is partial and dimensional,
and we describe it as such.

We do not invoke any extension of equation (20) beyond the FGHMV
setting. The FGHMV result is established for asymptotically AdS
spacetimes with spherical entangling surfaces and vacuum-state
perturbations. Extensions to other settings (flat-space, non-vacuum
states, non-spherical regions, large perturbations) are areas of active
research [12,13], and we do not claim that the escrow framework's
identification (20) carries through to those extensions automatically.
The framework's modular-Hamiltonian content as expressed in (20) is,
like the FGHMV result, presently a statement about the linear-response /
small-perturbation regime in asymptotically AdS-like settings.


\section{Consolidated Guardrails}

Sections II.E, III.F, and IV.E each conclude with explicit statements of
what the escrow framework does not commit to in the local context of
that section. We consolidate those statements here as a single list, for
use in subsequent work and in defense against the failure modes
documented in the framework's history.


\subsection{No escrow stress-energy tensor}

The escrow framework does not introduce a separate "escrow
stress-energy tensor" T\textsubscript{μν}\textsuperscript{escrow} to be added to the matter
stress-energy on the right side of Einstein's equation. The reason is
structural and matches the discussion in §II.E: Einstein's equations
are an algebraic identity between the matter stress-energy tensor
T\textsubscript{μν}\textsuperscript{matter} and the geometric Einstein tensor G\textsubscript{μν}, and the Bianchi
identities (∇\textsuperscript{μ} G\textsubscript{μν} = 0) require ∇\textsuperscript{μ} T\textsubscript{μν}\textsuperscript{matter} = 0 --- the matter
content already saturates the structure of conserved sources.
Gravitational binding energy is not a localizable tensorial quantity in
general relativity; it is encoded non-locally in the geometry itself
(the difference, for an isolated system, between the ADM mass at
infinity and the integrated rest mass on a spatial slice), and it cannot
be added to T\textsubscript{μν}\textsuperscript{matter} as an independent source. Attempting to add a
separate escrow source on the right side of Einstein's equation would
either violate the Bianchi identities or double-count the gravitational
configuration: the geometry on the left side already responds to the
full configuration through its non-local encoding of binding energy, and
the matter content on the right side already saturates the conservation
structure. The escrow framework supplies an interpretive dictionary for
the existing fields; it does not augment the field equation with any new
contribution. Jacobson's derivation [3] makes this constraint
particularly sharp: the entire content of $\delta Q$ on the right side of his
Clausius identity is already the matter stress-energy flux
T\textsubscript{μνk}\textsuperscript{μk}\textsuperscript{νdA} dλ, with no remaining room for an additional escrow
source. Attempts to extend the framework with additional source terms
reduce to either trivial restatements of the matter source (no new
physics) or violations of Jacobson's derivation (which would require
modifying the input $\delta Q$ = T·dS, not merely adding to the output).


\subsection{No vacuum computational substrate}

The escrow framework does not invoke a "vacuum computational
bandwidth," "information processing rate," or digital substrate
reading of the static escrow postulate. Such readings appear in some
popularizations of entropic gravity --- typically expressed as "the
vacuum is a computer with finite processing capacity, and gravity
emerges because the computer's bandwidth is partially consumed by
holding the escrow" --- but they are not part of the postulate as
stated in [1]. The local rate of escrow production at radius r is a
thermal quantity, set by $T_{U}$(r), and is correctly described in the
vocabulary of equilibrium thermodynamics. No additional ontological
structure is required to make the framework's predictions hold. The
framework's interpretation of why time dilation gives the rate it gives
(equation 10) is the thermodynamic Tolman-temperature reading, not a
computational-bandwidth reading.


\subsection{No claim about microscopic vacuum constituents}

The escrow framework, like Jacobson's derivation, treats the Einstein
equation as "an" equation of state of the vacuum rather than asserting
that spacetime is literally a thermodynamic system in some fundamental
sense. The thermodynamic analogy is structural --- the same Clausius
relation, the same area-as-entropy identification, the same local
equilibrium condition --- not ontological. We do not claim that
spacetime is built out of more fundamental thermodynamic degrees of
freedom in the way that an ideal gas is built out of molecules. Such a
claim would require identifying the microscopic constituents whose
statistical mechanics produces the area law; this is the project of
quantum gravity and is not addressed here. The escrow framework operates
entirely at the level of the macroscopic thermodynamic identity; it
makes no commitment to the underlying microphysics.


\subsection{Observer-dependence of $S_{\mathrm{esc}}$}

The escrow entropy $S_{\mathrm{esc}}$ is not observer-independent. The Unruh
temperature depends on the choice of accelerated observer; the local
Rindler horizon depends on the choice of local Rindler frame; the escrow
entropy inherits both dependences. What is observer-independent is the
relation $\delta Q$ = $T_{U}$ · d$S_{\mathrm{esc}}$ and its consequence (Einstein's equation,
equation 14), which hold for every observer in their own frame. This
observer-dependence is a feature of the thermodynamic interpretation,
not a bug: it is what distinguishes the escrow framework's reading of
GR from a global-entropy interpretation in which a unique S is assigned
to a spacetime. Attempts to define a global, observer-independent $S_{\mathrm{esc}}$
reduce either to the global entropy of the universe (which is not what
the static escrow postulate is about) or to the integrated horizon-flux
entropy (which is Jacobson's input, not new content).


\subsection{No empirical confirmation of equation (20)}

The lattice paper [2]'s partial recovery of the Bisognano--Wichmann
linear asymptote in 1+1D at small d1 with a 1/30 prefactor is not a
confirmation of equation (20). The 1/30 suppression is an unexplained
empirical observation; resolving its origin from first principles (UV/IR
mismatch, Dirichlet boundary correction, lattice regularization
artifact, or fundamental modification of the framework) is the principal
open theoretical question. We disclaim any reading of [2] as providing
empirical confirmation of the escrow framework's modular-Hamiltonian
identification, and we similarly disclaim any reading of [2]'s 3+1D
non-recovery as fundamentally ruling out the framework. The empirical
situation is partial, dimensional, and unresolved; we describe it as
such.


\subsection{No claim of new equations of motion}

The escrow framework does not produce new equations of motion. All
gravitational dynamics in the framework's purview reduce to Einstein's
equation (14), derived in the standard way via Jacobson [3] from $\delta Q$ =
T·dS, with the escrow framework's contribution being the interpretation
of the S as $S_{\mathrm{esc}}$. Attempts to extract new equations of motion from the
escrow postulate --- such as the C8 entropy-current attempt documented
in [9] --- have reduced to either rate-equation reformulations of
known results (no new physics) or to violations of the framework's
interpretive constraints (no new defensible physics). The framework's
mathematical content is the static escrow postulate (6) and the
consequences derived in [1]; its interpretive content is the
dictionary developed in the present paper. No new dynamical content is
claimed.


\subsection{$|U_{\mathrm{grav}}|$ is not given a single covariant definition}

Across the four legs of the translation, the symbol $|U_{\mathrm{grav}}|$ denotes
different quantities. In §II (test mass in a static gravitational
potential), $|U_{\mathrm{grav}}|$ = GMm/r is the Newtonian gravitational potential
energy of a probe particle. In §III (Schwarzschild horizon), $|U_{\mathrm{grav}}|$
= Mc²/2 is the self-binding energy of the configuration at horizon
scale, defined via GM²/$r_{s}$ with $r_{s}$ the Schwarzschild radius. In §III.D
(Reissner--Nordström and Kerr), $|U_{\mathrm{grav}}|$ is identified with the
chemical-potential decomposition (M − ΦQ − 2Ω\textsubscript{H} J)/2 that the Smarr
formula already supplies. In §IV (modular-Hamiltonian content of a
localized perturbation), $|U_{\mathrm{grav}}|$ is the energy of the localized
perturbation, formally identified with mc² for a point mass of rest mass
m. These four definitions are not connected by a single covariant,
observer-independent construction in standard general relativity.
Quasilocal energy proposals (Brown--York, Hawking mass, Komar integrals
for stationary configurations) are candidate constructions that could in
principle supply such a covariant definition, but the escrow framework
does not currently adopt one. We are explicit that the framework's
$|U_{\mathrm{grav}}|$ is presently a family of regime-specific identifications
glued together by common notation, and that providing a single covariant
definition is an open theoretical task. This limitation does not
undermine the consistency checks within each regime, but it does mean
that the framework's claim to a 'single thermodynamic object' is
partly notational: the object is $S_{\mathrm{esc}}$ in each regime, but the
$|U_{\mathrm{grav}}|$ feeding into $S_{\mathrm{esc}}$ is regime-specific.


\subsection{$T_{U}$ is used with two related but distinct conventions}

The Unruh temperature $T_{U}$ is used in this paper with two conventions
that are related by the Tolman shift but are not identical. In §II (test
mass in a Newtonian potential), $T_{U}$ is the local Unruh temperature of a
stationary observer at radius r, computed from the local proper
acceleration a = GM/r² via

$T_{U}$ = $\hbar$a/(2πc$k_{B}$) (local Unruh; §II convention)

This $T_{U}$ is what a local detector would register. In §III (Schwarzschild
horizon), $T_{U}$ is taken to equal the Hawking temperature

T\textsubscript{H} = $\hbar$κ/(2πc$k_{B}$), κ = c⁴/(4GM) (asymptotic Hawking; §III convention)

where κ is the surface gravity of the horizon. The Hawking temperature
is the redshifted (asymptotic) Unruh temperature of an observer hovering
just outside the horizon; the local Unruh temperature diverges at the
horizon, but its product with the Tolman redshift factor approaches the
finite asymptotic value T\textsubscript{H.} The two $T_{U}$ conventions are therefore
related by the Tolman shift

\[ T_{\mathrm{local}} \cdot  \sqrt{-g_{\mathrm{tt}}} = T_∞ \]

(the same exact Tolman relation invoked in §II.A), and they agree in the
appropriate limits. We use '$T_{U}$' uniformly in the body for notational
economy, with the understanding that the local vs. asymptotic
distinction is contextually determined by the regime under discussion.
This unification at the level of notation is not a unification at the
level of distinct physical observables: the local Unruh temperature and
the asymptotic Hawking temperature are different operational quantities,
related by the Tolman relation rather than identical.


\section{Conclusion}

The static gravitational entropy escrow postulate of [1] admits an
interpretive translation onto the established thermodynamic-gravity
literature. The translation comprises four legs. In Section II,
gravitational time dilation is recovered as the Tolman temperature shift
expressed in escrow language, with the dimensionless escrow of a test
mass identified as the escrow count

\[ 𝒩_{\mathrm{esc}} \equiv  2πr/λ_C \]

--- the number of reduced Compton wavelengths fitting into the orbit's
circumference. In Section III, Jacobson's derivation of Einstein's
field equations from $\delta Q$ = T·dS at local Rindler horizons becomes the
differential form of the static escrow postulate, with the
Bekenstein--Hawking entropy of a Schwarzschild black hole recovered
exactly (to all displayed digits, with no fudge factor) as

\[ S_{\mathrm{BH}}/k_B = 𝒩_{\mathrm{esc}}^{\mathrm{Schw}} = 4π(M/m_P)² \]

the escrow entropy of the configuration computed from its self-binding
energy and the Hawking temperature. In Section IV, the FGHMV
identification of the entanglement-entropy first law with linearized
Einstein equations becomes the statement that the escrow entropy of a
localized mass perturbation equals the modular-Hamiltonian content of
the wedge containing it,

\[ ΔK = 𝒩_{\mathrm{esc}}(d1)\qquad(\text{BW prediction}) \]

Section V consolidates the explicit guardrails against the failure modes
that the framework's history has surfaced.

None of these results modifies any standard equation of general
relativity, of quantum field theory, or of horizon thermodynamics. The
escrow framework's contribution at the level treated in this paper is
interpretive: it identifies a single thermodynamic object --- $S_{\mathrm{esc}}$ ---
through which four otherwise-separately-derived results are expressed as
faces of one identity. We note again (§V.G) that the 'single object'
description is partly notational, as the $|U_{\mathrm{grav}}|$ feeding into $S_{\mathrm{esc}}$
takes regime-specific forms across §II, §III, §III.D, and §IV that have
not yet been connected by a single covariant construction. The
framework's broader empirical anchors live in [1] and depend on the
horizon-thermodynamic content of the postulate, which is unchanged by
the present paper.

Three open theoretical problems are sharpened by the present
translation. First, the calculable origin of the 1/30 prefactor that
suppresses the lattice paper's 1+1D modular-Hamiltonian content
relative to the literal Bisognano--Wichmann asymptote remains
unresolved. A first-principles calculation of the lattice-regularized
modular Hamiltonian on the Dirichlet half-chain, accounting for the
UV/IR mismatch between localized lattice masses and the long-wavelength
BW prediction, is the calculation that would either explain the
suppression (rehabilitating the framework's modular-Hamiltonian
interpretation) or fail to explain it (constraining the framework). This
is the single most consequential follow-up calculation the framework
presently faces, and we identify it as the principal short-term research
priority.

Second, the 3+1D non-recovery of the BW asymptote within the lattice
sizes accessible to [2] is unresolved as to whether it reflects a
genuinely shorter characteristic length scale in higher dimensions or a
slower rate of approach to the continuum BW asymptote due to the smaller
conformal-symmetry group of the 3+1D massless free scalar. Lattice
computations at N ≳ 20 (feasible on dedicated GPU hardware) would
distinguish these interpretations and would constitute a decisive
next-generation test. A 2+1D test would supply a natural intermediate
dimensional case.

Third, the dimensional-dependence theorem connecting 1+1D and 3+1D
modular-Hamiltonian content is unidentified. Whether such a theorem
exists --- and what its form would be --- is a question that the FGHMV
framework, the Bisognano--Wichmann formula, and the lattice computations
of [2] are all positioned to address jointly. The escrow framework's
contribution to this question is the explicit identification of the
dimensionless quantity $\mathcal{N}_{\mathrm{esc}}$ ≡ 2πr/$\lambda_{C}$ as the candidate organizing
variable; whether it is the right variable, and whether its dimensional
dependence is calculable, are open.

A structural risk worth naming explicitly is that of retrofitted
inevitability: the possibility that the escrow framework's apparent
unification of the four results translated in this paper is an artifact
of having seen those four results first and then constructed a
dimensional combination ($S_{\mathrm{esc}}$ = $|U_{\mathrm{grav}}|$/$T_{U}$) that connects them in
retrospect, rather than the proposal of an organizing variable that
would have constrained the four results in advance had it been known.
The framework's escape from this failure mode is not interpretive but
calculational. Equation (20) of §IV makes a sharp numerical prediction
--- the modular-Hamiltonian content of a localized mass perturbation
should equal $\mathcal{N}_{\mathrm{esc}}$(d1) = 2π·m·d1 in natural units, with prefactor unity
--- and the lattice paper [2] tests it. The 1+1D result is a partial
confirmation (correct linear scaling, suppressed prefactor); the 3+1D
result is a non-recovery within accessible d1. Whether equation (20) is
a real identity with a calculable correction factor, or an arbitrary
dimensional combination that the lattice partially fits in 1+1D and
fails to fit in 3+1D for unrelated reasons, depends on whether the 1/30
prefactor has a first-principles origin in lattice-regularization
theory. A successful first-principles calculation would move the
framework from "interpretive synthesis with retrofitted-inevitability
risk" to "organizing variable with a calculable correction factor." A
failed calculation would constrain the framework's modular-Hamiltonian
leg. Either outcome is informative, and either outcome distinguishes the
framework's content from a pure relabeling exercise. The framework's
status as a candidate proposal rather than a derived theory is therefore
reducible to a single open calculational question, which we name as
such.

To make the falsification structure explicit, we identify five
conditions under which specific legs of the framework would be ruled
out. (i) If a first-principles lattice calculation accounting for
Dirichlet boundary conditions, lattice UV cutoff effects, and the finite
half-chain depth produces a prefactor that disagrees with the
lattice-measured 0.030--0.037 at the same protocol, the
modular-Hamiltonian leg expressed in equation (20) is constrained or
ruled out in its present form. (ii) If a 3+1D lattice computation at N ≥
24 with d1 swept up to d1 ≤ N/3 ≈ 8 fails to show approach toward the BW
linear asymptote (i.e., the peak-then-decay structure of ΔK(d1) observed
at N = 14 persists with the peak location not migrating significantly
toward larger d1 and the decay exponent not flattening toward zero), the
modular-Hamiltonian leg of the framework is constrained in 3+1D, the
3+1D non-recovery becomes a structural feature requiring explanation
rather than a finite-size artifact, and the framework retreats to its
horizon-thermodynamic and Tolman-redshift content only. We name N ≥ 24
with d1 ≤ N/3 as the concrete threshold rather than leaving 'any
accessible d1' unbounded; computation at larger N remains informative
but the framework would already be constrained at this threshold. (iii)
If a charged or rotating black hole calculation reveals that the escrow
postulate cannot recover $S_{\mathrm{BH}}$ for a configuration where the first law of
black hole mechanics applies (i.e., a case not absorbed by the Smarr
formula's chemical-potential decomposition), the horizon-thermodynamic
leg of the framework is constrained. (iv) If a strong-field calculation
(e.g., the exact Schwarzschild redshift for a compact object, or the
analogous Kerr or Reissner--Nordström redshift) shows that the escrow
translation of equation (9) produces a quantitatively wrong prediction
when extended beyond the Newtonian weak-field limit --- i.e., the
iterated escrow identity does not reproduce √(−$g_{tt}$) at strong field
within reasonable precision --- the Tolman-redshift leg is constrained
to remain a weak-field effective description, and the framework
explicitly disclaims any strong-field reading of §II's translation. (v)
If a rigorous theorem shows that Jacobson's $S_{\mathrm{horizon}}$ cannot be
identified with any quasilocal energy construction (such as Brown--York
energy, Hawking mass, or a Komar integral on the horizon) in a manner
consistent with the escrow postulate's structural form S = |U|/T ---
i.e., if the open theoretical task named in §V.G is shown to have no
positive resolution --- the framework's claim that $S_{\mathrm{esc}}$ and $S_{\mathrm{horizon}}$
refer to the same underlying object is undermined, and the framework's
Jacobson-leg content reduces to a notational coincidence rather than a
structural identification. In each case the falsification is partial,
removing one leg of the five-condition translation without necessarily
ruling out the others. We note these explicitly rather than relying on
implicit logical structure.

We commit, in advance of the calculations named in conditions (i)--(v),
to the following pre-registered retractions. If condition (i) is
satisfied (the first-principles prefactor calculation disagrees with the
lattice-measured 0.030--0.037 at the same protocol), the
modular-Hamiltonian leg of the framework will be withdrawn, and any
successor version of this paper will publish only the time-dilation
(§II) and Jacobson/Bekenstein--Hawking (§III) legs. If condition (ii) is
satisfied (the N ≥ 24 lattice persistently shows the 3+1D
peak-then-decay structure), the 3+1D applicability of the framework will
be withdrawn; the modular-Hamiltonian leg will be scoped explicitly to
1+1D only. If condition (iii) is satisfied (a Smarr-resistant black-hole
configuration fails to recover $S_{\mathrm{BH}}$), the horizon-thermodynamic leg will
be withdrawn. If condition (iv) is satisfied (the strong-field test
produces a quantitatively wrong redshift prediction in the iterated
escrow translation), §II will be explicitly scoped to the Newtonian
weak-field limit and no strong-field claim will be defended. If
condition (v) is satisfied (the §V.G open task is shown to have no
positive resolution), the finite-level identification $S_{\mathrm{horizon}}$ = $S_{\mathrm{esc}}$
for static black holes (§III.C and §III.D) will be retracted as a
structural claim and retained only as a notational coincidence; the
differential-level identification δ$S_{\mathrm{horizon}}$ = δ$S_{\mathrm{esc}}$ (§III.B), which is
the content of the escrow reading of Jacobson's derivation and depends
on the first law rather than on a covariant binding-energy construction,
does not depend on the §V.G open task and would remain intact. These
pre-registered commitments are intended to make the framework's
falsification structure costly rather than rhetorical: a Popperian
theory carries with it the obligation to specify, in advance, what would
constitute its failure and what one would do in response. We accept that
obligation here and bind any future version of this paper to honor it.

The escrow framework joins an existing body of thermodynamic and
entropic interpretations of gravity. Jacobson's 1995 derivation [3],
to which Section III is most directly indebted, derives Einstein's
equations from $\delta Q$ = T·dS at local Rindler horizons. Verlinde's 2011
entropic gravity proposal and its 2017 emergent-gravity extension take a
related but distinct approach, deriving Newton's law from an entropy
gradient on a holographic screen and relating the deep-MOND scale to the
de Sitter horizon temperature. Padmanabhan's emergent-gravity program
has long emphasized that gravitational field equations have the standing
of a thermodynamic equation of state and that horizon thermodynamics is
the natural variable for understanding gravity. The covariant entropy
bound (Bousso) constrains the entropy of light-sheets and provides a
related but operationally distinct entropy/area relation. The escrow
framework's specific contribution within this literature is the
identification of the dimensionless ratio $\mathcal{N}_{\mathrm{esc}}$ ≡ 2πr/$\lambda_{C}$ (and its
horizon analog $\mathcal{N}_{\mathrm{esc}}^{\mathrm{Schw}}$ ≡ 4π(M/$m_{P}$)²) as a candidate organizing
variable for the thermodynamic interpretation, and the explicit
modular-Hamiltonian identification of equation (20) as the falsifiable
empirical content. We do not claim that the escrow framework supersedes
Jacobson, Verlinde, Padmanabhan, or Bousso, or that it occupies a unique
position relative to them; we claim that it is a specific proposal
within the broader thermodynamic-gravity program, with the conceptual
translation of this paper showing how it connects to the established
literature and with the lattice paper [2] supplying its principal
falsifiable test. For transparency on the citation chain: references
[3]--[8] and [10]--[13] are peer-reviewed publications, while
references [1], [2], and [9] are the present author's Zenodo
preprints that have not yet been submitted to peer review; the
framework's empirical anchor (the lattice paper [2]) and its
theoretical anchor (the framework paper [1]) are therefore the
author's own work, awaiting independent replication and peer review.

The framework as it presently stands is best characterized as a
candidate organizing variable for the thermodynamic interpretation of
gravity, with the static escrow postulate as the proposed organizing
relation, the lattice paper as the empirical anchor for its
modular-Hamiltonian interpretation, the framework paper as the empirical
anchor for its horizon-thermodynamic interpretation, and the present
paper as the explicit conceptual re-reading connecting these to the
established thermodynamic-gravity literature. The translation is
internally consistent and recovers standard results without
modification. Whether the candidate organizing variable proves to be the
right one is a question that the open theoretical problems above will
determine, and that future empirical and analytical work will
adjudicate.


\subsection*{Data, Code \& Resources}
\noindent Manuscript source and archived record: \url{https://github.com/Windstorm-Institute/escrow-spacetime} (Zenodo: \href{https://doi.org/10.5281/zenodo.20126090}{10.5281/zenodo.20126090}). Experiments, datasets, and analysis code: \url{https://github.com/Windstorm-Labs/escrow-spacetime}. Windstorm Institute: \url{https://windstorminstitute.org}. Windstorm Labs: \url{https://windstormlabs.com}.

\section*{References}
\small

\textbf{[1]} G. L. Whitmer III, Gravitational Entropy Escrow: An
Interpretive Synthesis of Thermodynamic Approaches to Gravity, Zenodo
(2026), doi:10.5281/zenodo.20031931.

\textbf{[2]} G. L. Whitmer III, A Lattice QFT Test of
the Static Escrow Postulate: 1+1D and 3+1D Falsification, with a
Suggestive Modular-Hamiltonian Window in 1+1D, Zenodo (2026),
doi:10.5281/zenodo.20057537.

\textbf{[3]} T. Jacobson, Thermodynamics of Spacetime: The Einstein
Equation of State, Phys. Rev. Lett. 75, 1260 (1995),
arXiv:gr-qc/9504004.

\textbf{[4]} R. C. Tolman and P. Ehrenfest, Temperature Equilibrium in a
Static Gravitational Field, Phys. Rev. 36, 1791 (1930).

\textbf{[5]} T. Faulkner, M. Guica, T. Hartman, R. C. Myers, and M. Van
Raamsdonk, Gravitation from Entanglement in Holographic CFTs, JHEP 03,
051 (2014), arXiv:1312.7856.

\textbf{[6]} N. Lashkari, M. B. McDermott, and M. Van Raamsdonk,
Gravitational Dynamics from Entanglement "Thermodynamics", JHEP 04,
195 (2014), arXiv:1308.3716.

\textbf{[7]} J. D. Bekenstein, Black Holes and Entropy, Phys. Rev. D 7,
2333 (1973).

\textbf{[8]} S. W. Hawking, Particle Creation by Black Holes, Commun.
Math. Phys. 43, 199 (1975).

\textbf{[9]} G. L. Whitmer III, On the Status of Local Entropy-Current
Extensions of the Gravitational Entropy Escrow Framework: A
Clarification Note, Zenodo (2026), doi:10.5281/zenodo.20041991.

\textbf{[10]} S. Ryu and T. Takayanagi, Holographic Derivation of
Entanglement Entropy from AdS/CFT, Phys. Rev. Lett. 96, 181602 (2006),
arXiv:hep-th/0603001.

\textbf{[11]} J. J. Bisognano and E. H. Wichmann, On the Duality
Condition for Quantum Fields, J. Math. Phys. 17, 303 (1976).

\textbf{[12]} H. Casini, M. Huerta, and R. C. Myers, Towards a derivation
of holographic entanglement entropy, JHEP 05, 036 (2011),
arXiv:1102.0440.

\textbf{[13]} T. Faulkner, F. M. Haehl, E. Hijano, O. Parrikar, C.
Rabideau, and M. Van Raamsdonk, Nonlinear Gravity from Entanglement in
Conformal Field Theories, JHEP 08, 057 (2017), arXiv:1705.03026.


\end{document}

